% ================================
% === Reviewer 5 ===
% ================================
\reviewerblock{Reviewer~5}{%
We thank the reviewer for the positive assessment of the decoupling perspective and the module design, and for the precise minor corrections.}

\cresp{Figure~2 caption inconsistency: The caption refers to ``(Top) UTD'' and ``(Bottom) SLR-C'', but the figure shows ``A'' as the top and ``B'' as the bottom. Please revise the caption or the figure to ensure consistent correspondence.}
{Thank you for identifying this inconsistency. To satisfy the page limit, we shorten the Figure~2 caption and remove the former ``Top/Bottom'' panel description. The revised caption now summarizes the functions of SLR-C and UTD without making the inconsistent correspondence.
\revloc{the rewritten panel-independent caption of Fig.~\ref{fig:fdil_framework}}
}
\showFdilFramework

\cresp{Missing experimental settings in Table~5:: Two additional configurations should be included: (i) Lexical + Scenario, and (ii) Canonical + Scenario.}
{We add the two rows to the prior-source ablation in Table~\ref{tab:prior_ablation} and rerun them under the same $K{=}10$ reranking-only protocol as every existing row. We also revise the paragraph below the table so that it compares all three pairwise combinations: both new pairs improve over their single-source components, Canonical+Scenario is the strongest pair on Hard-F1, and the three-source ensemble remains best on aggregate Avg-F1. Thus, the revision changes both the table coverage and the interpretation rather than simply inserting two numbers.
\revloc{the added Lexical+Scenario and Canonical+Scenario rows in Table~\ref{tab:prior_ablation}; the revised comparison paragraph immediately following the table in Heterogeneous Priors in SLR-C}
}
\showPriorAblation

\cresp{Unexplained anomaly in Table~6: Provide a justification for why $K=5$ $K=28$ yield better performance than $K=10$.}
{We revisit this anomaly by replacing the former single-seed comparison with a three-seed sweep of $K\in\{5,10,28\}$ under the final protocol. Table~\ref{tab:k_sensitivity_test} now reports mean$\pm$std for the complete metric set and identifies \res{$K=10$} as the best aggregate setting, so the earlier single-seed dip is no longer used in our conclusion. We rewrite the accompanying interpretation to explain both sides of the trade-off: $K=5$ can exclude plausible positive labels, whereas $K=28$ admits all classes and introduces distractors without improving semantic coverage; $K=10$ best balances candidate recall and local separability. We also change the default setting stated in the Method and Implementation Details to $K=10$ for consistency.
\revloc{the replaced three-seed results and mean$\pm$std entries in Table~\ref{tab:k_sensitivity_test}; the rewritten trade-off paragraph in Sensitivity to Candidate Size $K$; the consistent $K=10$ setting in Semantic Local Reranking and Implementation Details}
}
\showKSensitivity

\cresp{Typo in Appendix tables: In Table~B.1 and Table~B.2 (page~35), ``FDIR'' should be corrected to ``FDIL''.}
{Thank you for identifying this typo. During revision, the former Appendix Tables~B.1--B.2 are removed because their validation-split results were redundant with the revised main analysis, and the manuscript had to remain within 35 pages; consequently, the two reported ``FDIR'' instances are no longer present. We nevertheless perform a manuscript-wide search across the text, captions, and tables, and correct or confirm every remaining acronym occurrence as ``FDIL.''
\revloc{removal of the former Appendix Tables~B.1--B.2; acronym consistency checked throughout the main text, table entries, and captions}
}
